% A "piece of advice" note for a practitioner forecasting rideshare demand
% at 1-6 week horizons. Distilled from a study of week-ahead highway-traffic
% forecasting (which shares the same skeleton).
% Compile: pdflatex rideshare_forecasting_advice.tex
\documentclass[11pt]{article}
\usepackage[margin=1in]{geometry}
\usepackage{amsmath,amssymb}
\usepackage{enumitem}
\usepackage{xcolor}
\usepackage[colorlinks=true,linkcolor=blue,urlcolor=blue]{hyperref}
\setlength{\emergencystretch}{3em}
\setlist[enumerate]{itemsep=4pt,topsep=4pt}

\title{\textbf{Forecasting Rideshare Demand at 1--6 Week Horizons}\\[3pt]
\large A few pieces of advice}
\author{}
\date{}

\begin{document}
\maketitle
\vspace{-2.5em}

\noindent\hrulefill

\paragraph{Where this comes from.} These notes are distilled from a detailed study
of \emph{week-ahead highway-traffic} forecasting. Rideshare demand shares its
skeleton --- a dominant weekly\,$\times$\,daily seasonal cycle driven by human
schedules, strong holiday/event sensitivity --- so the structural lessons transfer.
But rideshare differs from highway traffic in ways that \emph{change where modeling
pays}, and the 1--6 week horizon is its own regime. Here's what I'd tell you.

\section{Start here (and don't skip it)}

\begin{enumerate}[leftmargin=*]
\item \textbf{Build the calendar-adjusted seasonal baseline first, and make it the
bar everything must clear.} Per zone, the rule is:
\begin{center}
\fbox{\begin{minipage}{0.9\linewidth}\centering\vspace{2pt}
\textbf{forecast} $=$ (recency-weighted average of the same day-of-week-and-hour
over the last few weeks) $\times$ (event/holiday factor on special days)
\vspace{2pt}\end{minipage}}
\end{center}
The seasonal average uses geometrically-decaying weekly weights (one parameter);
the factor is the historical ratio (actual\,/\,baseline) on that holiday/event,
pooled over past occurrences. On highway flow this reached $R^2\approx0.95$ and a
weekly-aggregate error around $3.5\%$, and \emph{nothing fancier beat it} ---
autoregression, gradient boosting, pooled panel models, cross-unit features all
tied or lost. Expect a similarly brutal baseline for rideshare. Anything you deploy
must beat \emph{this}, in a leakage-free rolling backtest, before it earns its
complexity.

\item \textbf{Know which ``band'' 1--6 weeks lives in.} Write demand as
\emph{level $\times$ seasonal shape $+$ noise}. The forecast horizon is a low-pass
filter: as it grows, the short-range, spatial, surge-propagation signal (what
real-time dispatch models exploit) decays away. By 1--6 weeks you are forecasting
the \emph{seasonal backbone plus a slowly-moving level} --- almost none of the
minutes-ahead structure survives. Practical consequence: \textbf{do not reach for
the spatio-temporal deep nets built for short-horizon dispatch}; they shine at
minutes, not weeks. Different horizon, different toolbox.
\end{enumerate}

\section{Where a learned model genuinely helps at 1--6 weeks}

The honest framing: at this horizon a model cannot extract hidden signal from
demand's \emph{own past} --- the seasonal baseline already captures it. Its value is
in \emph{learning a known response} and \emph{reducing variance}. Rideshare,
happily, has far more of both than highway traffic.

\begin{enumerate}[leftmargin=*,resume]
\item \textbf{Your \#1 lever: learn the known-ahead event \& calendar response.}
Rideshare is much more event-driven than traffic --- concerts, sports, conventions,
festivals, holidays, New Year's Eve, airport surges. The killer property is that
these are \textbf{known weeks ahead}, so they need no covariate forecast. Generalize
the hand-crafted holiday factor into a \emph{learned} event-response model
(gradient boosting or hierarchical) with event metadata (venue, capacity, type)
interacted with zone and time-of-day, plus per-zone heterogeneity. This is the
clearest place a learned model beats the plain baseline at this horizon.

\item \textbf{Model the level with macro / known covariates --- it is more dynamic
than traffic's.} Highway traffic is a near-utility with a near-frozen level;
rideshare is a \emph{discretionary} service whose level moves with employment,
tourism (flight arrivals, hotel occupancy), fuel prices, competitor presence,
marketing, and your own adoption curve. That level is partly forecastable from
slowly-varying or known-ahead covariates. A structural / state-space model
(level $+$ trend $+$ seasonal $+$ regression), or a global model with these inputs,
is the right tool --- and it matters \emph{increasingly toward 6 weeks}, where level
staleness, not the seasonal shape, dominates the error.

\item \textbf{Exploit your panel --- pool across zones and cities.} You almost
certainly have far more cross-sectional units than a sensor network (zones, hexes,
cities) plus coupled series (demand, supply, ETA). A single \emph{global} model
(DeepAR / temporal-fusion-transformer / global gradient boosting) with zone
embeddings and shared covariates pools strength: it denoises sparse, low-demand
zones and learns the event/calendar response far better than per-zone fits. This is
where ``more data'' genuinely pays --- concentrated on the \emph{noisy-zone tail} and
the \emph{shared response function}, not on mining each zone's own history.

\item \textbf{Forecast distributions, not just point means.} Supply positioning,
driver incentives, and surge decisions care about the \emph{spread and the tails},
not the average. Even where your point forecast only ties the baseline, calibrated
quantiles --- especially event-driven and regime-dependent uncertainty --- are
high-value and genuinely model-amenable. Train on quantile/pinball loss and report
calibration, not just MAPE.
\end{enumerate}

\section{Traps that will quietly cost you}

\begin{enumerate}[leftmargin=*,resume]
\item \textbf{Respect the known-ahead constraint --- and the surge-leakage trap.} At
1--6 weeks you may only use covariates whose \emph{future} value is known at
inference. Events, holidays, calendar, scheduled promotions and planned price
changes: \emph{yes}. \textbf{Weather: no} --- you cannot forecast it 1--6 weeks out;
at best fold in climatology, and never train on weather \emph{actuals} you won't
have in production. \textbf{Price / surge is endogenous}: if you train demand on
realized surge you leak badly (you won't know future surge, and surge itself
responds to demand). Model demand at a reference price, or treat \emph{your own
planned} pricing as a known input. This leakage silently inflates backtests and
breaks on deployment.

\item \textbf{Assume the seasonal shape drifts.} Traffic's weekly shape is frozen
year to year; rideshare's is not (adoption, behavior change, work-from-home shifting
commute\,$\leftrightarrow$\,leisure mix). Use \emph{adaptive} seasonality (EWMA over
weeks, or a time-varying seasonal), refit on a schedule, and monitor for drift and
regime breaks (a pandemic, a city launch, a competitor's exit). A static seasonal
profile will slowly rot.

\item \textbf{Get the evaluation right --- this is where most effort is wasted or
misled.}
  \begin{itemize}[itemsep=2pt,topsep=2pt]
  \item Proper \emph{rolling-origin} backtests, with the known-ahead constraint
    enforced (parameters fit only on data prior to each origin).
  \item Beware \textbf{MAPE on low-demand zones}: tiny denominators explode it (we
    saw $>$500\% on quiet hours that were forecast fine in absolute terms).
    Aggregate, or use weighted / scaled / quantile metrics.
  \item \textbf{Bucket} the errors --- event vs.\ normal days, by zone density, by
    regime. The headline number hides exactly where you win and lose.
  \item Always report \textbf{skill over the baseline} (item 1), not raw error.
  \item Estimate the \textbf{predictability ceiling} (e.g.\ an oracle in-sample fit)
    so you know how much headroom exists before you chase it.
  \end{itemize}
\end{enumerate}

\section{The mindset}

\begin{enumerate}[leftmargin=*,resume]
\item \textbf{Complexity should follow information, not precede it.} A transformer
will not beat a calendar average on the deterministic backbone --- that part is
solved. Spend model complexity only where rideshare carries exploitable information
\emph{beyond} the regular pattern: events, the macro-driven level, panel pooling,
and the decision-relevant tails. The honest sequence is always: \emph{nail the
baseline $\to$ measure where the variance and the headroom actually are $\to$ point
the learned model exactly there.}
\end{enumerate}

\vspace{0.5em}
\noindent\hrulefill

\paragraph{The good news.} Highway traffic was the extreme case --- almost all
deterministic backbone, little else --- which is why a one-parameter rule was
unbeatable. Rideshare moves a lot of mass into the components where models earn
their keep: events, macro coupling, a large panel, regime change, and a
decision-coupled objective. So a well-aimed learned model has \emph{real} room here.
The whole trick is that it be \emph{aimed}, not merely deployed.

\end{document}
